\sectionkicker{Theme synthesis}
\subsection*{横断比較 — 長時間Agentでは状態をどこで持つかが設計になる}
\addcontentsline{toc}{subsection}{横断比較 — 長時間Agentでは状態をどこで持つかが設計になる}
SwarmX、ChatGPT Dreaming、SmoothAgentは同じagent研究・機能ではないが、長時間実行を一回のmodel callとして扱えない点で共通する。scheduler、persistent memory、context transformationという別々の状態管理層へ分けると関係が見えやすい。
\medskip
\begin{center}
\small
\renewcommand{\arraystretch}{1.16}
\begin{tabularx}{\linewidth}{@{}p{0.20\linewidth}X@{}}
\toprule
観察軸 & 6月の一次資料・論文から読めること \\
\midrule
実行スケジューリング & SwarmXはmulti-call / tool-using agent workloadをschedulerの対象として扱い、単一inference requestではなく依存を持つ実行列としてserving問題を定義する。 \cite{src-e8d9fd84bddc} \\
persistent memory & ChatGPT Dreamingはreviewable memory summaryとstaged rolloutを伴うmemory architectureとして説明され、長期contextを単純にpromptへ積み続ける以外のproduct-level state managementを示した。 \cite{src-c3ab33c62226} \\
context transformation & SmoothAgentはlong-horizon servingでcontext transformationを先回りして実行するsystems researchであり、contextそのものをruntime側で加工・準備する設計を扱う。報告性能はauthor claimの範囲に留める。 \cite{src-80f8b4d899ff} \\
model intelligenceとの境界 & 3件を横断すると、長時間Agentの品質はmodel capabilityだけでなく、いつ実行するか、何を記憶するか、どのcontextを渡すかというruntime/state設計にも依存する。各資料は異なる層を扱うため性能値を直接比較しない。 \cite{src-e8d9fd84bddc,src-c3ab33c62226,src-80f8b4d899ff} \\
\bottomrule
\end{tabularx}
\renewcommand{\arraystretch}{1.0}
\end{center}
\normalsize
\begin{claimboundary}[この比較の境界]
この表は本号で既に検証・参照している一次資料と論文を横断比較の形へ再配置したものである。vendor / project / author が報告した性能値や優位性は独立再現済みの一般則へ変換せず、本文とSource Notesの帰属境界をそのまま維持する。
\end{claimboundary}
